\section{Introduction}\label{sec:introduction}

Urban heat is produced by interactions among surface materials, built form,
vegetation, moisture, radiation, ventilation, and human activity. The Local
Climate Zone (LCZ) framework made this complexity operational by defining a
standard climate-form vocabulary for urban temperature studies
\citep{stewart2012LocalClimate}. Earth-observation classification and the
Worldwide Urban Database and Access Portal Tools (WUDAPT) subsequently made
LCZ mapping repeatable across cities \citep{bechtel2012ClassifiLocal,
bechtel2015MappingLocal,ching2018WudaptUrban,bechtel2019GeneratiWudapt}, and a
a global LCZ product supports comparative urban and Earth-system analysis
\citep{demuzere2022GlobalLocal}. The framework is therefore indispensable for
comparability, but comparability is not the same as local planning resolution.

Within-class thermal variability has been observed across climates, cities,
seasons, and measurement systems \citep{geletic2016LandSurface,
skarbit2017EmployinUrban,kwok2019WellDoes,yang2020InvestigDiversit,
li2022VariabilLand}. This matters especially in coastal metropolitan regions,
where sea proximity, topographic transition, vegetation structure, and
contrasting compact, open, industrial, and peripheral fabrics may produce
substantial variation within a nominal LCZ. Recent work consequently relates
LST to two- and three-dimensional morphology, vegetation, and functional
context rather than treating class membership as a complete explanation
\citep{zhou2022UnderstaEffects,wang2023EffectivFactors,
zhou2023IdentifyEffects,liu2025UnveilinDifferen,yu2026UrbanMorpholo}. The
question is no longer only which LCZ occupies a location, but which measurable
features make locally similar labels climatically different.

Three methodological gaps remain. First, many classification pipelines learn
LCZ labels from multisource imagery but inherit label uncertainty and spatial
scale assumptions \citep{qiu2018FeatureImportan,yoo2019ComparisBetween,
qiu2020MultilevFeature,yoo2020ImprovinLocal,yan2022ComparinObject}. Second,
urban climate models can be accurate while remaining difficult to translate
into design action if 3D structure, street interface, vegetation volume, and
network context are compressed into opaque predictors. Third, random train-test
splits can reward spatial autocorrelation rather than transferable morphology;
grid and block comparisons show that analytical support changes estimated
form--temperature relationships \citep{jeon2023ImpactsUrban,wu2022ScaleEffect,
li2026InvestigEffect}. Explainability, spatial holdouts, and scale sensitivity
are therefore parts of the estimand, not optional presentation devices.

Izmir is a demanding test of this proposition. Its functional urban region
contains dense coastal districts, hillside settlement, open low-rise fabrics,
large-footprint industrial zones, suburban expansion, and fragmented
peripheral development. Coastal heat-risk interpretation must also account for
local vulnerability and the limits of transferring LCZ semantics between
cities \citep{canguzel2024ClimateChange,zhang2025ApplicabLocal}. We therefore
treat the global LCZ map as weak-label infrastructure and combine it with a
local 5 m surface model, exact building and road geometry, vegetation structure, remotely
sensed land cover and LST, coastal context, and competition-adjusted
network accessibility.

The study makes four evidence-backed contributions. First, it documents an
analysis-ready 250 m Izmir grid with explicit provenance, missingness,
collinearity, and proxy language. Second, it quantifies the incremental value
of morphology, surface-model structure, vegetation, green-blue context, and street-network
features under leakage-aware spatial cross-validation. Third, it replaces
Euclidean green-space buffers with a network-cost two-step floating catchment
area (2SFCA) measure and tests 400, 800, and 1,200 m thresholds. Fourth, it
connects weak-label morphology to LST, feature importance, uncertainty, and
leave-district-out transfer. This explanatory sequence follows the broader
principle that urban-form machine learning should connect prediction,
interpretation, and planning meaning \citep{karl2026RelationBetween}.

The analysis combines a prespecified 569-location manual audit with weak-label and tabular model evaluation. The audit provides a direct comparison between global labels, model predictions, and locally reviewed morphology, so the reported profiles and predictions are interpreted against observed local fabric.
